\documentclass[11pt]{article}

\usepackage[T1]{fontenc}
\usepackage[margin=1in]{geometry}
\usepackage{lmodern}
\usepackage{microtype}
\usepackage{amsmath,amssymb,mathtools}
\usepackage{amsthm}
\usepackage{booktabs}
\usepackage{array}
\usepackage{multirow}
\usepackage{xcolor}
\usepackage{enumitem}
\usepackage{needspace}
\usepackage{tikz}
\usetikzlibrary{arrows.meta,positioning,fit}
\usepackage[hidelinks]{hyperref}
\usepackage[nameinlink,noabbrev]{cleveref}
\usepackage{listings}
\usepackage{url}
\usepackage[backend=biber,style=numeric,sorting=nyt,maxbibnames=99]{biblatex}
\addbibresource{references.bib}

\newtheorem{theorem}{Theorem}
\newtheorem{lemma}[theorem]{Lemma}
\newtheorem{proposition}[theorem]{Proposition}

\definecolor{deepblue}{RGB}{23,69,124}
\definecolor{darkred}{RGB}{145,32,32}
\definecolor{lightgray}{RGB}{245,246,248}
\ifdefined\ANONYMOUS
  \newcommand{\pdfpaperauthor}{Anonymous}
\else
  \newcommand{\pdfpaperauthor}{Justin Thaler, Kai Zhang, Scott Kominers, Quang Dao}
\fi
\hypersetup{
  colorlinks=true,
  linkcolor=deepblue,
  citecolor=deepblue,
  urlcolor=deepblue,
  pdftitle={Symmetry-Quotient Forgery Attacks on Odd-Characteristic SNOVA},
  pdfauthor={\pdfpaperauthor},
  pdfsubject={Cryptanalysis of the public odd-characteristic SNOVA drafts and previews},
  pdfkeywords={SNOVA, multivariate signatures, forgery attack, NIST PQC}
}
\setlist{nosep,leftmargin=*}
\lstset{basicstyle=\ttfamily\small,backgroundcolor=\color{lightgray},frame=single,breaklines=true}

\newcommand{\F}{\mathbb F}
\newcommand{\Mat}{\operatorname{Mat}}
\newcommand{\Sym}{\operatorname{Sym}}
\newcommand{\rank}{\operatorname{rank}}
\newcommand{\Span}{\operatorname{span}}
\newcommand{\MQ}{\mathsf{MQ}}
\newcommand{\Hashimoto}{\mathsf{Hashimoto}}
\newcommand{\transpose}{\mathsf T}
\newcommand{\SNOVA}{\textsf{SNOVA}}
\newcommand{\QRUOV}{\textsf{QR-UOV}}
\newcommand{\MAYO}{\textsf{MAYO}}
\newcommand{\bits}{\text{ bits}}
\newcommand{\attacktitle}{Symmetry-Quotient Forgery Attacks on Odd-Characteristic \texorpdfstring{\SNOVA}{SNOVA}}

\title{\attacktitle}

\ifdefined\ANONYMOUS
  \author{Anonymous submission}
\else
  \author{%
    Justin Thaler\thanks{a16z crypto research and Georgetown University} \and
    Kai Zhang\thanks{MIT} \and
    Scott Kominers\thanks{Harvard University and a16z crypto research} \and
    Quang Dao\thanks{Carnegie Mellon University and LayerZero Labs}%
  }
\fi
\date{}

\begin{document}
\setlength{\emergencystretch}{2em}
\raggedbottom
\maketitle

\ifdefined\DISCLOSURE
\begin{center}
\fcolorbox{darkred}{white}{\parbox{0.91\textwidth}{\centering\small
\textbf{Confidential coordinated disclosure.} This draft is intended for the \SNOVA{} submission team and NIST. Please do not redistribute before the coordinated public-disclosure date.}}
\end{center}
\fi

\begin{abstract}
We give a forgery attack on the odd-characteristic form of \SNOVA{}, a third-round candidate in NIST's Additional Digital Signatures process.  \SNOVA{}'s original characteristic-two parameters have already been broken, and the team's public repair moves to odd characteristic, where the public-key matrices become symmetric; we show that this repair also falls below its targeted security.  At Levels~I, III, and V our estimates drop to about $2^{138.9}$, $2^{184.3}$, and $2^{228.0}$ Boolean gates, against targets $2^{143}$, $2^{207}$, and $2^{272}$.  All nine Version~2.3 draft parameter sets fall short, by 4.06 to 44.05 bits; the six Version~2.4 analysis-preview shapes (conditional on their retaining the symmetric construction, since no matching specification is currently public) fall short as well, albeit by as little as 1.43 bits.

The attack builds on the existing Beullens forgery framework (Eurocrypt 2025), which recasts forgery as solving a system of quadratic equations whose hardness is governed by the rank of an associated linear map.  We show that the new symmetry forces that map to depend only on the symmetric combinations of its inputs, so the antisymmetric combinations are invisible to the attack and the rank drops from $m_1\ell^2$ to $m_1\binom{\ell+1}{2}$ (here $\ell$ is the block size and $m_1$ a scheme parameter).  This leaves a much smaller system to solve: across the nine Version~2.3 sets, from 48 quadratic equations in 112 variables to 96 in 224.

The reduction is exact and holds for every key with symmetric public matrices.  We also confirm it directly against \SNOVA{}'s official Level-I reference implementation, reconstructing the reduced system from its own key material.

All bit-security figures follow the same cost model used in the \SNOVA{} team's own analysis; we do not claim a full-size forgery.  The attack does not apply to the unsymmetrized $q=16$ public matrices.
\end{abstract}

\section{Introduction}\label{sec:intro}

\SNOVA{} is a structured Oil-and-Vinegar signature scheme designed to combine very small public keys with short signatures and fast verification.  NIST advanced \SNOVA{} to the third round of the Additional Digital Signatures process in May~2026, while allowing the surviving candidates to submit further specifications and implementations~\cite{NISTIR8610,NISTRound3}.  The scheme has not, however, had a single stable public parameter generation, and no definitive Round~3 parameter package was public when this manuscript was prepared.

The NIST-posted Round~2 package uses $\F_{16}$ and nonsymmetrized public base matrices.  That line has already been heavily damaged: Beullens reduced the security margins by exploiting the structured whipping layer~\cite{Beullens2024SNOVA}, and the subsequent structure-aware wedge attack of Bros et al.\ reports that all Round~2 parameters are broken, while noting a remaining rank conjecture in its complexity analysis~\cite{BrosEtAl2026SNOVAWedge}.  The \SNOVA{} team's later public drafts pursue odd-characteristic variants of the scheme intended to evade these attacks.  Version~2.1 added prime-field variants with symmetric public matrices; Version~2.3 supplies $q=19$ implementations and KATs; the public analysis repository later posted a six-set ``Version~2.4'' preview; and the authors' reformulation paper presents odd characteristic as a principal route for redesigning the scheme~\cite{SNOVA23,SNOVA24Preview,DingEtAl2026Reformulating}.  An odd-characteristic attack therefore targets the public repair strategy rather than the already-vulnerable submitted parameters.

The Version~2.3 attack chapter retains a forgery analysis in the style of Beullens~\cite{Beullens2024SNOVA}.  There, forgery is recast as solving a system of quadratic equations, and the size of that system is taken to be $m_1\ell^2$: the analysis lets the powers $S^0,\dots,S^{\ell-1}$ of the public matrix $S$ contribute all $\ell^2$ ordered pairs $(a,b)$, and it models the resulting equations as essentially random of that size.  The draft itself states that this chapter is inherited from Version~2.0 and still needs updating~\cite{SNOVA23}.

This paper shows that the ordered model is incompatible with the new public symmetry.  Put $A=\F_q[S]$.  The irreducibility of the characteristic polynomial makes $A$ a field isomorphic to $\F_{q^\ell}$, while $S=S^\transpose$ makes every element of $A$ symmetric.  For fixed $u$ and public base form $P_i$, the Gram pairing
\[
 (\xi,\eta)\longmapsto
 u^\transpose(I_n\otimes\xi)P_i(I_n\otimes\eta)u
\]
is a symmetric $\F_q$-bilinear form on the $\ell$-dimensional space $A$.  It therefore factors through $\Sym^2_{\F_q}(A)$.  Over $\F_{19}$, the ordered tensor square decomposes as
\[
 A\otimes_{\F_q}A
 =\Sym^2_{\F_q}(A)\oplus\bigwedge\nolimits^2_{\F_q}(A),
\]
and the entire alternating summand is invisible to quadratic evaluation.  In coordinates this is the equality of the $(a,b)$ and $(b,a)$ features.  The feature dimension is therefore
\[
 K=m_1\binom{\ell+1}{2},
\]
not $m_1\ell^2$.  This identity is universal for every key generated with symmetric public matrices; it is not a rare MinRank event and cannot be removed by adding more $A,B,Q$ terms (the ``emulsifier'' matrices of the whipping layer, which randomize the public map).

\subsection{Results}

Recall the parameters: $v$ vinegar and $o$ oil matrix variables (so $n=v+o$), block size $\ell$, and $r$ columns per signature block, with $m_1=\lceil or/\ell\rceil$ scalar-expanded public base matrices and $M=or\ell$ scalar verification equations.  Our main algebraic result is the following.

\Needspace{19\baselineskip}
\begin{theorem}[Universal symmetry quotient, informal]\label{thm:intro}
Consider a \SNOVA{} public key for which the power matrix $S$ and the scalar-expanded public base matrices $P_1,\ldots,P_{m_1}$ are symmetric.  After any affine-column substitution in the Beullens framework, the homogeneous quadratic part factors through
\[
 \bigoplus_{i=1}^{m_1}\Sym^2_{\F_q}(\F_q[S]).
\]
Consequently its rank is at most $K=m_1\binom{\ell+1}{2}$.  When both the quotient rank and the affine-constraint rank are generic---$\rho=K$ and $\lambda_R=M-K$, the case realized in every relation we tested---forging a selected target reduces directly to solving $K$ quadratic equations in
\[
 N_{\rm res}=\ell(v+o)-\left(or\ell-K\right)
\]
variables.
\end{theorem}

The full theorem, including nongeneric ranks, appears as \Cref{thm:rankbound,thm:residual}.  The rank cap comes from the domain of the quadratic feature map.  The $A,B,Q$ emulsifier matrices may randomize the induced map on the symmetric quotient, but they cannot recover the alternating coordinates discarded before emulsification.  The generic full-rank case is the one realized in every relation we tested (\Cref{sec:validation}).

Applying the submission's own MQ methodology gives the Version~2.3 estimates in \Cref{tab:main}.  The submission's estimator admits two conventions for the semi-regular series, differing by whether one homogenizing variable is inserted.  The public analysis script and the values printed in Version~2.3 Table~11 use the convention without that variable ($p_1=0$); inserting the extra variable ($p_1=1$) yields a stricter, higher estimate.  We adopt the stricter $p_1=1$ convention for the headline figures, so our reported costs are conservative relative to the submission's own numbers; even so, every $q=19$ set falls below target.  Applying the identical quotient to the Version~2.4 analysis preview gives the additional estimates in \Cref{tab:snova24}.

\begin{table}[t]
\centering
\caption{Symmetry-reduced residual systems and estimated attack costs.  ``Old'' is the Beullens-forgery estimate in \SNOVA{} Version~2.3 Table~11.  ``New'' uses the stricter $p_1=1$ convention we adopt for our headline figures.  Negative margins mean that the estimate is below the category target.}\label{tab:main}
\small
\setlength{\tabcolsep}{4.2pt}
\begin{tabular}{@{}c l c c c r r@{}}
\toprule
Level & Parameters $(v,o,\ell,r)$ & Residual MQ & Old & New & Target & Margin\\
\midrule
I   & $(28,5,4,4)$   & $50$ in $102$ & $194$ & $138.94$ & $143$ & $-4.06$\\
I   & $(48,16,2,2)$  & $48$ in $112$ & $164$ & $130.30$ & $143$ & $-12.70$\\
I   & $(28,4,4,5)$   & $50$ in $98$  & $194$ & $138.94$ & $143$ & $-4.06$\\
\midrule
III & $(40,7,4,4)$   & $70$ in $146$ & $264$ & $184.29$ & $207$ & $-22.71$\\
III & $(72,24,2,2)$  & $72$ in $168$ & $234$ & $184.29$ & $207$ & $-22.71$\\
III & $(38,5,4,5)$   & $70$ in $142$ & $239$ & $184.29$ & $207$ & $-22.71$\\
\midrule
V   & $(50,9,4,4)$   & $90$ in $182$ & $333$ & $227.95$ & $272$ & $-44.05$\\
V   & $(96,32,2,2)$  & $96$ in $224$ & $303$ & $238.77$ & $272$ & $-33.23$\\
V   & $(52,6,4,6)$   & $90$ in $178$ & $333$ & $227.95$ & $272$ & $-44.05$\\
\bottomrule
\end{tabular}
\end{table}

\SNOVA{}'s verifier rejects any signature whose blocks are symmetric.  Because the attack is built from symmetric column relations, one might hope this rule blocks it.  We show it does not: in the square variants the attacker adds fixed offsets that force a nonzero antisymmetric entry into every block, so the forged signature passes the check while the system to be solved is unchanged, and in the rectangular variants the blocks are not symmetric objects in the first place.

We also confirm the reduction against the reference code.  Reconstructing the official Level-I implementation reproduces the reduced system exactly from its shipped key material, and fixing one signature column in fact leaves a nearly square system of 50 quadratics in 52 variables (\Cref{sec:crossnormal,sec:validation}).  A direct search for a shortcut beyond solving that system---a shared solution locus or an exploitable low-rank structure---found none.

\subsection{Organization}

\Cref{sec:related} surveys prior work.  \Cref{sec:prelim} fixes notation and recalls the affine-column framework.  \Cref{sec:symmetry} proves the quotient and its multibase Gram extension.  \Cref{sec:attack} derives the smaller system of equations a forger must solve, and \Cref{sec:rejection} shows that the verifier's rule against symmetric signatures does not block the attack.  \Cref{sec:crossnormal} gives the cross-column 50-in-52 normal form.  \Cref{sec:complexity} evaluates our attack's runtime against both the Version~2.3 parameter sets and the Version~2.4 preview.  \Cref{sec:validation} reports the experiments confirming the reduction against \SNOVA{}'s reference implementation and testing its solving-degree model.  \Cref{sec:countermeasures} gives countermeasures.  \Cref{app:noAshortcut} rules out an obvious alternative attack---a MinRank shortcut over the extension field---and \Cref{app:multibase} develops the multibase solver extension.

\section{Related Work}\label{sec:related}

\paragraph{Beullens's forgery attack.}
Beullens identified \SNOVA{}'s structured whipping layer, expressed the public map through ordered bilinear power-pair coordinates, and showed that low-rank combinations of the corresponding emulsifier matrices lead to forgery attacks~\cite{Beullens2024SNOVA}.  The Version~2.3 countermeasure increases the number and diversity of $A,B,Q$ terms to make the induced ordered-coordinate matrix resemble a random matrix.

Our attack uses the same affine-column architecture but a different source of rank loss.  We do not find a low-rank element in a family of otherwise full-rank ordered-coordinate matrices.  We show that the odd-characteristic homogeneous feature vector itself lies in a $K$-dimensional symmetric subspace.  Thus the quotient applies to every key and every relation, even when the induced map is maximally random and full rank on that quotient.

\paragraph{Stable-ideal attacks.}
Cabarcas et al. exploit the action of the commutative group generated by $S$ to design a polynomial-system solver specialized to \SNOVA{}~\cite{Cabarcas2025SNOVA}.  Their work can be viewed as improving the residual-system solving stage.  The present result instead reduces the number of residual equations before solving begins.  The approaches are compatible, but we do not combine their cost benefits here.

\paragraph{Characteristic-two wedge attacks and the odd-characteristic redesign.}
The structure-aware wedge attack exploits \SNOVA{}'s block-ring structure and reports that the NIST Round~2 $q=16$ parameters are broken, while explicitly retaining a conjectural rank component in the analysis~\cite{BrosEtAl2026SNOVAWedge}.  The \SNOVA{} authors subsequently investigated the wedge map and proposed a reformulation emphasizing odd-characteristic and symmetric-form variants~\cite{TsengEtAl2026Wedge,DingEtAl2026Reformulating}.  Our attack is complementary: it does not apply to the unsymmetrized $q=16$ public matrices, but it attacks the public odd-characteristic repair direction through an elementary quotient that is exact and independent of the wedge-rank conjecture.

\paragraph{Key-recovery attacks.}
Several works interpret the scalar expansion of \SNOVA{} as UOV, or lift it to smaller extension-field UOV instances, and transfer key-recovery attacks~\cite{IkematsuAkiyama2024,LiDing2024,NakamuraEtAl2025}.  Our result is a universal forgery reduction and does not recover the secret oil space.

\paragraph{Relation to the symmetric-algebra key-recovery attack.}
The decomposition
$V\otimes V=\Sym^2(V)\oplus\bigwedge^2(V)$ and the vanishing of the alternating
part under quadratic evaluation are classical; we claim no novelty for that
algebra~\cite{Lam2005}.  The closest cryptanalytic precedent is the
$p$-reduced symmetric-algebra framework of Jin, Pan, He, Gong, and
Ding~\cite{JinEtAl2026}, which generalizes Ran's characteristic-two wedge
attack~\cite{Ran2026} to odd characteristic.  Three differences are worth
stating.  First, theirs is a \emph{key-recovery} XL attack that reconstructs the
oil space, whereas ours is a universal forgery that recovers no secret.  Second,
their algebra is the quotient $\F_q[x_1,\dots,x_n]/(x_1^p,\dots,x_n^p)$ on all
$n=v+o$ variables, in which the polar form is a degree-two element and the
Macaulay computation runs to degree up to $o(p-1)$; our symmetric square is the
degree-two part of the $\ell$-dimensional power algebra $\F_q[S]$, and collapses
only the ordered power pairs of the Beullens feature map.  Third, their cost is
strongly characteristic-sensitive---it grows with $p$ and is ineffective at the
large characteristics of \QRUOV{}---whereas our quotient dimension
$m_1\binom{\ell+1}{2}$ is independent of $p$ for odd $p$.  The two attacks are
therefore complementary, and ours applies directly in the moderate-characteristic
regime $q=19$ that the odd-characteristic \SNOVA{} drafts adopt.  Their own
concrete recommended-parameter evaluation targets \QRUOV{} and reports no
improvement over the existing reconciliation attack at its large characteristic;
they provide no \SNOVA{} cost tables, so no numerical head-to-head on the $q=19$
Version~2.3 sets is available.  Association-scheme or Delsarte machinery concerns
related objects but is not the mechanism here.

\paragraph{Concurrent structural cryptanalyses of UOV-derived signatures.}
This paper is part of a concurrent set of self-contained structural
cryptanalyses of UOV-derived signature schemes~\cite{PolarHodgeMAYO,FasterQRUOV}.
The other works study polar-flag quotients and Hodge contractions for \MAYO{}
and \MAYO{}-style constructions, and extension-field slicing and pseudo-oil
attacks on \QRUOV{}.  The attacks are technically independent, but share the
broader lesson that the algebra used to compress a UOV construction must be
incorporated explicitly into its security analysis.  The closest of these to the
present work concerns \MAYO{}: it identifies preimage-preserving polar quotients
of the characteristic-two \MAYO{} verification map~\cite{PolarHodgeMAYO}.  Both
results expose a deterministic loss of attack-relevant dimension before generic
solving, so that added downstream randomization cannot restore the discarded
directions.  The constructions are nonetheless different: the \MAYO{} quotient
acts on the input and output of the quadratic map and comes with a lifting
theorem, whereas the present quotient acts on \SNOVA{}'s ordered power-pair
feature space through a symmetric-square factorization, with no analogous
lifting step.

\paragraph{Our contribution.}
The contribution is recognizing that the public odd-characteristic redesign makes this elementary quotient load-bearing in the existing \SNOVA{}
affine-column attack; deriving the exact rank and residual-system consequences
for square and rectangular variants; giving the structural
$\Sym^2(A)$ proof and its multibase Gram extension; proving that the advertised
signature rejection does not block the attack; and validating the reduction
against the current implementation and cost methodology.  We found no
corresponding symmetry quotient in the Version~2.3 attack chapter or the Version~2.4 analysis preview.  The former continues to use ordered $\ell^2$ coordinates and explicitly states that it still requires updating.

\paragraph{Negative transfer to generic outer tensors.}
The \SNOVA{} reduction depends on a deterministic symmetry of the public feature map.  In tests of tensor-UOV constructions with independent random outer tensors, two-column restrictions already saturated the complete output space and yielded no complementary affine equations.  Thus the cross-column quotient is not a generic consequence of UOV compression; it is tied to the concrete symmetry introduced by the odd-characteristic \SNOVA{} design.

\section{Preliminaries and the Existing Forgery Framework}\label{sec:prelim}

\subsection{Parameters and scalar expansion}

Let $q$ be a prime power, let $\ell$ be the size of the square matrix blocks, and let $r$ be the number of columns of each signature block.  The public map has $n=v+o$ matrix variables.  Version~2.3 defines
\[
 m_1=\left\lceil\frac{or}{\ell}\right\rceil
 \quad\text{and}\quad
 M=or\ell,
\]
where $m_1$ is the number of scalar-expanded public base matrices and $M$ is the number of scalar output equations.  A signature is an $n$-tuple of blocks in $\Mat_{\ell\times r}(\F_q)$.  Equivalently, writing its columns as
\[
 U=[u_0\mid u_1\mid\cdots\mid u_{r-1}],
 \qquad u_j\in\F_q^{n\ell},
\]
the attack starts with $n\ell$ variables after expressing the other columns affinely in terms of $u_0$.

Let $S\in\Mat_{\ell}(\F_q)$ be the public symmetric matrix with irreducible characteristic polynomial, and put
\[
 \Sigma_a=I_n\otimes S^a\in\Mat_{n\ell}(\F_q),
 \qquad 0\le a<\ell.
\]
For each base matrix $P_i\in\Mat_{n\ell}(\F_q)$ define the bilinear power-pair forms
\begin{equation}\label{eq:B}
 B_i^{(a,b)}(x,y)=x^\transpose\Sigma_aP_i\Sigma_b y.
\end{equation}
The public odd-characteristic algorithms generate each scalar-expanded public matrix symmetrically: diagonal blocks are symmetric and opposite off-diagonal blocks are transposes.  Thus $P_i^\transpose=P_i$ for the public $q=19$ drafts, which we confirm directly against the official Level-I KAT in \Cref{sec:validation}.

\subsection{The Beullens affine-column substitution}

Beullens observed that the \SNOVA{} public map can be written as a linear transformation of the values $B_i^{(a,b)}(u_j,u_k)$~\cite{Beullens2024SNOVA}.  The attacker selects matrices $R_j\in\F_q[S]$ and offsets $v_j\in\F_q^{n\ell}$, with $R_0=I$, and substitutes
\begin{equation}\label{eq:relation}
 u_j=(I_n\otimes R_j)u+v_j,
 \qquad 0\le j<r.
\end{equation}

after identifying $u=u_0$ and taking $v_0=0$ if desired.

For a fixed relation $R=(R_0,\ldots,R_{r-1})$, the substituted verification map for a target $t\in\F_q^M$ can be written
\begin{equation}\label{eq:affine-system}
 E_R z_{\rm ord}(u)+L_Ru+c_R=t.
\end{equation}
Here
\[
 z_{\rm ord}(u)
 =\left(B_i^{(a,b)}(u,u)\right)_{1\le i\le m_1,\ 0\le a,b<\ell}
 \in\F_q^{m_1\ell^2},
\]
$E_R\in\F_q^{M\times m_1\ell^2}$ depends only on public constants and the chosen column relation, while $L_Ru+c_R$ collects all cross and constant terms involving the offsets.

In the ordered-coordinate analysis, if $\rank(E_R)=\rho$, Gaussian elimination supplies $M-\rho$ affine-linear equations and leaves at most $\rho$ independent quadratic equations.  The Version~2.3 attack chapter models $E_R$ as essentially random and bases its concrete estimate on $\rho$ near the ordered-coordinate maximum.  Our observation is that, whenever the public base matrices are symmetric, $E_R$ is applied to a feature vector with deterministic coordinate equalities.

\section{The Symmetry Quotient}\label{sec:symmetry}

\begin{lemma}[Power-pair symmetry]\label{lem:powerpair}
Suppose $S^\transpose=S$ and $P_i^\transpose=P_i$.  Then for every $u\in\F_q^{n\ell}$ and every $0\le a,b<\ell$,
\[
 B_i^{(a,b)}(u,u)=B_i^{(b,a)}(u,u).
\]
\end{lemma}

\begin{proof}
The quantity is a scalar, so it equals its transpose:
\[
 \begin{aligned}
 u^\transpose\Sigma_aP_i\Sigma_bu
 &=\left(u^\transpose\Sigma_aP_i\Sigma_bu\right)^\transpose\\
 &=u^\transpose\Sigma_b^\transpose P_i^\transpose\Sigma_a^\transpose u
 =u^\transpose\Sigma_bP_i\Sigma_au.
 \end{aligned}
\]
We used that every power of the symmetric matrix $S$ is symmetric.
\end{proof}

\begin{theorem}[Structural symmetric-square factorization]\label{thm:structural-sym2}
Let
\[
 A=\F_q[S]\cong \F_{q^\ell}
\]
be viewed as an $\ell$-dimensional vector space over $\F_q$.  For a fixed
$u\in\F_q^{n\ell}$ and public base matrix $P_i$, define
\[
 \Phi_{u,i}(\xi,\eta)
 =u^\transpose(I_n\otimes\xi)P_i(I_n\otimes\eta)u,
 \qquad \xi,\eta\in A.
\]
Then $\Phi_{u,i}$ is a symmetric $\F_q$-bilinear form on $A$.  Consequently
it factors uniquely through $\Sym^2_{\F_q}(A)$, and the combined homogeneous
feature map for $P_1,\ldots,P_{m_1}$ factors through
\[
 \bigoplus_{i=1}^{m_1}\Sym^2_{\F_q}(A).
\]
When the characteristic is odd, the killed subspace in the ordered tensor
model is canonically
\[
 \bigoplus_{i=1}^{m_1}\bigwedge\nolimits^2_{\F_q}(A),
\]
of dimension $m_1\binom{\ell}{2}$.
\end{theorem}

\begin{proof}
The irreducibility of the characteristic polynomial makes $A$ a field, and
$S=S^\transpose$ makes every element of $A$ symmetric.  Bilinearity over
$\F_q$ is immediate.  For $\xi,\eta\in A$, the value is a scalar and hence
is equal to its transpose:
\[
\begin{aligned}
 \Phi_{u,i}(\eta,\xi)
 &=u^\transpose(I_n\otimes\eta)P_i(I_n\otimes\xi)u\\
 &=\left(u^\transpose(I_n\otimes\eta)P_i(I_n\otimes\xi)u\right)^\transpose\\
 &=u^\transpose(I_n\otimes\xi)P_i(I_n\otimes\eta)u
 =\Phi_{u,i}(\xi,\eta).
\end{aligned}
\]
Thus $\Phi_{u,i}$ is symmetric and factors through the universal symmetric
square.  In odd characteristic,
$A\otimes_{\F_q}A=\Sym^2_{\F_q}(A)\oplus\bigwedge^2_{\F_q}(A)$, so the
alternating summand is exactly the discarded part.
\end{proof}

Concretely, for $\ell=4$, the ordered power
space has dimension $16$, while its symmetric square has dimension $10$;
the six missing directions are exactly $\bigwedge^2(A)$.  With $m_1=5$,
this gives the 30 affine-linear output directions observed in the Level-I
end-to-end experiment.

Define the unordered feature vector
\[
 z_{\rm sym}(u)
 =\left(B_i^{(a,b)}(u,u)\right)_{1\le i\le m_1,\ 0\le a\le b<\ell}
 \in\F_q^K,
 \qquad K=m_1\binom{\ell+1}{2}.
\]
Let
\[
 D:\F_q^K\longrightarrow\F_q^{m_1\ell^2}
\]
be the duplication map that sends an unordered coordinate $(i,\{a,b\})$ to both ordered positions $(i,a,b)$ and $(i,b,a)$ when $a\ne b$, and to the single diagonal position when $a=b$.  Then
\begin{equation}\label{eq:dup}
 z_{\rm ord}(u)=Dz_{\rm sym}(u).
\end{equation}

\begin{figure}[t]
\centering
\resizebox{\textwidth}{!}{%
\begin{tikzpicture}[node distance=11mm,>=Latex,
 box/.style={draw,rounded corners,align=center,minimum height=10mm,minimum width=29mm,fill=lightgray}]
\node[box] (u) {$u\in\F_q^{n\ell}$};
\node[box,right=of u] (sym) {$z_{\rm sym}(u)$\\$K=m_1\binom{\ell+1}{2}$};
\node[box,right=of sym] (ord) {$z_{\rm ord}(u)$\\$m_1\ell^2$ positions};
\node[box,right=of ord] (out) {quadratic output\\in $\F_q^M$};
\draw[->] (u)--node[above]{quadratic}(sym);
\draw[->] (sym)--node[above]{$D$}(ord);
\draw[->] (ord)--node[above]{$E_R$}(out);
\draw[->,deepblue,thick,bend right=23] (sym) to node[below]{$\overline E_R=E_RD$} (out);
\end{tikzpicture}%
}
\caption{The ordered feature vector lies in the image of a $K$-dimensional duplication map.  Randomizing $E_R$ cannot increase the rank of the composite beyond $K$.}\label{fig:factor}
\end{figure}

\begin{theorem}[Symmetry rank bound]\label{thm:rankbound}
For every public key whose power matrix and scalar-expanded public base matrices are symmetric, and for every affine-column relation $R$, the following holds.  The homogeneous quadratic part in \Cref{eq:affine-system} factors as
\[
 \overline E_R z_{\rm sym}(u),
 \qquad \overline E_R=E_RD\in\F_q^{M\times K}.
\]
In particular,
\[
 \rank(\overline E_R)\le K=m_1\binom{\ell+1}{2}.
\]
\end{theorem}

\begin{proof}
The structural factorization is \Cref{thm:structural-sym2}.  In the power
basis $1,S,\ldots,S^{\ell-1}$, it becomes the coordinate identity
\Cref{eq:dup}.  Substituting that identity into \Cref{eq:affine-system}
gives the displayed factorization, and the rank bound follows from the number
of columns of $\overline E_R$.
\end{proof}

\paragraph{Why the emulsifier does not remove the relation.}
The Version~2.3 changes use many independently generated $A_{i,\alpha},B_{i,\alpha},Q_{i,\alpha,1},Q_{i,\alpha,2}$ matrices so that the ordered-coordinate $E_R$ resembles a general matrix rather than the repeated block-diagonal matrix attacked in the first-round construction.  This may make $\overline E_R$ a generic map on its $K$-dimensional domain.  It cannot make that domain larger.  The lost antisymmetric coordinates vanish before the emulsifier is applied.

\paragraph{The reduction is universal, not a weak-key event.}
Prior analyses search over relations $R$ to find an unusually low rank.  Here the cap $K$ holds for every relation and every public key satisfying the public-matrix symmetry rule.  A further rank drop below $K$ can only improve the residual system relative to the full-rank quotient analyzed below.

\subsection{The Gram viewpoint and several base vectors}\label{sec:gram}

The same argument gives a clean generalization of the attack.  Choose
$c$ independent base vectors $x_1,\ldots,x_c$ and an $A$-valued linear
column relation
\begin{equation}\label{eq:cbase}
 u_j=\sum_{s=1}^{c}(I_n\otimes R_{j,s})x_s+v_j,
 \qquad R_{j,s}\in A.
\end{equation}
Put $V_c=\F_q^c\otimes_{\F_q}A$, of dimension $c\ell$.  For each public
base matrix $P_i$, the homogeneous restriction of the verifier is a symmetric
Gram pairing on $V_c$.  We therefore obtain the following exact extension.

\begin{proposition}[$c$-base Gram factorization]\label{prop:cbase}
Under \Cref{eq:cbase}, the homogeneous feature map factors through
\[
 \bigoplus_{i=1}^{m_1}\Sym^2_{\F_q}(V_c),
\]
and hence through at most
\[
 K_c=m_1\binom{c\ell+1}{2}
\]
quadratic coordinates.  If $c=r$ and the $r\times r$ relation matrix over
$A$ is invertible, the relation is merely a public change of basis among the
signature columns and imposes no restriction on the signature space.
\end{proposition}

\begin{proof}
Apply the proof of \Cref{thm:structural-sym2} to the labels
$e_s\otimes\xi\in V_c$.  Symmetry of the public form identifies a pair of
labels with its reversal, so the universal feature space is the symmetric
square of $V_c$.  An invertible $A$-linear column change is bijective.
\end{proof}

The Gram formulation also closes a broader search over arbitrary coefficient
matrices.  For a general relation with $R_{j,s}\in\Mat_\ell(\F_q)$, define
the right power-orbit space
\[
 \mathcal W_s=\Span_{\F_q}\{S^aR_{j,s}:0\le a<\ell,\ 0\le j<r\}.
\]
If base $s$ is active, some $R_{j,s}$ is nonzero.  The map
$A\to\Mat_\ell(\F_q)$, $\xi\mapsto\xi R_{j,s}$, is injective because every
nonzero $\xi\in A$ is invertible.  Therefore
\begin{equation}\label{eq:orbit-lower}
 \dim_{\F_q}\mathcal W_s\ge\ell.
\end{equation}
An $A$-valued relation attains equality.  Thus allowing arbitrary block
coefficients cannot use fewer than $\ell$ power directions per active base;
it normally makes the orbit space much larger.  In 500 trials per
$c\in\{1,2,3,4\}$ on the Level-I shape, $A$-valued relations always had
total orbit dimension $c\ell$, whereas unrestricted random
$\ell\times\ell$ coefficients always had dimension $c\ell^2$.  Accidental
rank defects of the public Gram map remain possible, but a broader coefficient
search cannot improve the universal feature count by shrinking the underlying
power-orbit space below the $A$-valued construction.

For the official Level-I constants, the quotient ranks for $c=1,2,3,4$ are
respectively $50,80,80,80$, equal to the generic maximum
$\min\{K_c,80\}$ in every tested relation.  The $c\ge2$ maps are therefore
surjective: they remove the salt-membership phase but leave a structured
nonlinear inversion problem.  We discuss the resulting candidate multibase
attack and its present solver uncertainty in \Cref{app:multibase}.

\section{The Symmetry-Assisted Forgery Attack}\label{sec:attack}

Let
\[
 \rho=\rank(\overline E_R)\le K.
\]
Choose a full-rank matrix $W_R\in\F_q^{(M-\rho)\times M}$ whose rows span the left kernel of $\overline E_R$.  Multiplying \Cref{eq:affine-system} by $W_R$ eliminates every quadratic term:
\begin{equation}\label{eq:linear}
 W_RL_Ru=W_R(t-c_R).
\end{equation}
Let $\lambda_R=\rank(W_RL_R)$.  Gaussian elimination expresses $\lambda_R$ coordinates of $u$ as affine functions of the remaining
\[
 N_R=n\ell-\lambda_R
\]
coordinates.  Substitution into a row basis of \Cref{eq:affine-system} leaves at most $\rho$ quadratic equations.

\begin{theorem}[Residual MQ reduction]\label{thm:residual}
Fix a relation and offsets, and suppose the linear system \Cref{eq:linear} is consistent for the target $t$.  Then forging $t$ reduces in polynomial time to solving at most $\rho$ quadratic equations in $n\ell-\lambda_R$ variables, where
\[
 \rho\le m_1\binom{\ell+1}{2},
 \qquad
 \lambda_R\le M-\rho.
\]
In the generic case $\rho=K$ and $\lambda_R=M-K$, the coefficient matrix $W_RL_R$ has full row rank $M-\rho$, so \Cref{eq:linear} is consistent for \emph{every} target, and the residual system consists of
\begin{equation}\label{eq:resparams}
 \boxed{
 K=m_1\binom{\ell+1}{2}\text{ quadratics in }
 N_{\rm res}=n\ell-M+K\text{ variables}.}
\end{equation}
A solution extends by linear algebra to a signature satisfying the original $M$ verification equations.
\end{theorem}

\begin{proof}
The left-kernel step gives \Cref{eq:linear}.  When that system is consistent, solve a maximal independent subsystem and substitute into the original equations.  Modulo the affine constraints, the quadratic coefficient vectors lie in the $\rho$-dimensional column space of $\overline E_R$, so a row basis gives at most $\rho$ independent quadratics.  Reversing the substitution produces $u$, and \Cref{eq:relation} gives all signature columns.  Because the reduction used row operations and exact substitutions, the resulting signature satisfies all $M$ original equations.  When $\rho=K$ and $\lambda_R=M-K$, the matrix $W_RL_R$ is $(M-\rho)\times n\ell$ of rank $M-\rho$, hence surjective, so \Cref{eq:linear} solves for any right-hand side.
\end{proof}

\begin{proposition}[Effect of a rank defect]
\label{prop:affine-rank}
Two departures from the generic ranks behave differently.
\begin{enumerate}
\item A \emph{quotient-rank} defect $\rho<K$ accompanied by full affine rank $\lambda_R=M-\rho$ leaves $\rho$ quadratic equations in $n\ell-M+\rho$ variables and imposes no condition on the target.  The gap $(n\ell-M+\rho)-\rho=n\ell-M$ is independent of $\rho$, so such a defect only decreases the number of quadratic equations and cannot make the attack harder.
\item An \emph{affine-rank} defect $\lambda_R<M-\rho$ makes \Cref{eq:linear} impose $(M-\rho)-\lambda_R$ consistency conditions on the target.  For a target drawn through the hash, these hold with probability about $q^{-((M-\rho)-\lambda_R)}$, so the attacker resamples the salt---or searches over relations and offsets for a full-rank affine subsystem---at that expected cost.  A defect of this kind is therefore \emph{not} free.
\end{enumerate}
\end{proposition}

\begin{proof}
(1) With $\lambda_R=M-\rho$, the matrix $W_RL_R$ has full row rank, so \Cref{eq:linear} is consistent for every target; the elimination pins $M-\rho$ coordinates and leaves $n\ell-M+\rho$ variables in $\rho$ quadratics.  (2) $W_RL_R$ has $M-\rho$ rows and rank $\lambda_R$, so its image is a $\lambda_R$-dimensional subspace of $\F_q^{M-\rho}$; a right-hand side $W_R(t-c_R)$ induced by a uniform target lies in that image with probability $q^{-((M-\rho)-\lambda_R)}$.
\end{proof}

\paragraph{The tabulated systems are the generic full-rank case.}
The residual dimensions in our tables are stated for $\rho=K$ and $\lambda_R=M-K$, where the affine subsystem has full row rank, \Cref{eq:linear} is unconditionally solvable, and the reduction is a genuine selected-target forgery.  This is the case realized in every relation we tested (\Cref{sec:validation}), where we verify both ranks directly.  A quotient-rank defect would only lower the equation count; an affine-rank defect would add the target-consistency cost of \Cref{prop:affine-rank}(2), which we do not observe.

\paragraph{Attack algorithm.}
The reduction is a universal forgery in the key-only setting: the attacker targets any message it chooses and makes no signing queries.  For an arbitrary message $\mathcal M$:
\begin{enumerate}
\item sample a salt and compute the target $t=H(\mathcal M\|\textsf{salt})$;
\item choose a public column relation $R$ and offsets $v_j$;
\item construct $\overline E_R,L_R,c_R$ and row-reduce as above;
\item solve the residual underdetermined MQ system;
\item reconstruct $u$ and then the complete signature matrix $U$ using \Cref{eq:relation}.
\end{enumerate}
The first, second, fourth, and fifth steps are the same attack architecture used by Beullens.  The new ingredient is that step~3 must quotient the homogeneous features by symmetry before estimating their rank.

\section{Bypassing the Symmetric-Block Rejection}\label{sec:rejection}

Version~2.3 notes that signatures composed of powers of $S$ can be dangerous when the public matrices are symmetric.  For odd $q$, the verifier therefore rejects square signature blocks that are symmetric: it allows none when $\ell>2$ and at most $\lfloor n/4\rfloor$ when $\ell=2$.

The affine-column attack has fixed offsets in addition to its free vector $u$.  Those offsets can enforce nonsymmetry identically.

\begin{lemma}[A constant-skew column relation]\label{lem:skew}
Assume $\ell\ge2$ and that the characteristic polynomial of $S$ is irreducible.  Let $e_0,e_1$ be two distinct standard basis vectors of $\F_q^\ell$.  There exists $R_1\in\F_q[S]$ such that
\[
 e_0^\transpose R_1=e_1^\transpose.
\]
With $R_0=I$, choose offsets for the first two signature columns so that, in each matrix block,
\[
 (v_1)_0-(v_0)_1=\delta
\]
for a fixed $\delta\in\F_q^\times$.  Then every assignment to $u$ satisfies
\[
 U[0,1]-U[1,0]=\delta
\]
in that block.  In particular, the block is never symmetric.
\end{lemma}

\begin{proof}
Because the characteristic polynomial is irreducible, $\F_q^\ell$ is a simple $\F_q[S]$-module.  Hence the nonzero row vector $e_0^\transpose$ is cyclic under right multiplication by $\F_q[S]$, and some $R_1$ maps it to $e_1^\transpose$.  Under the relation $u_0=u+v_0$ and $u_1=(I_n\otimes R_1)u+v_1$, the variable-dependent parts of the two entries are $e_1^\transpose u$ and $e_0^\transpose R_1u$, which coincide.  Their difference is therefore the chosen offset constant $\delta$.
\end{proof}

The construction can be applied independently to every block.  It does not change the homogeneous quadratic map $\overline E_R$; offsets affect only $L_R$ and $c_R$.  For the square $\ell=2$ sets it makes every block nonsymmetric, which is stronger than the verifier's threshold.  For the square $\ell=4$ sets it satisfies the zero-symmetric-block requirement.  The rectangular variants have blocks in $\Mat_{\ell\times r}$ with $r\ne\ell$, so matrix symmetry is not defined and no such rejection is applied.

\section{A Cross-Column Gram Normal Form}\label{sec:crossnormal}

The multibase Gram viewpoint yields a second exact reduction for the official Level-I square parameters.  Write the four signature columns as $x_0,x_1,x_2,x_3\in\F_{19}^{132}$.  The complete symmetry-reduced feature space consists of four self blocks and six cross blocks.  For a pair $s<t$, let
\[
 G_{st}(x_s,x_t)
 =\bigl(B_i^{(a,b)}(x_s,x_t)\bigr)_{1\le i\le5,\ 0\le a,b<4}
 \in\F_{19}^{80},
\]
and let $E_{st}:\F_{19}^{80}\to\F_{19}^{80}$ be the corresponding public output map.

\begin{proposition}[Official Level-I cross-pair certificate]\label{prop:crosspair}
For the official Version~2.3 Level-I $q=19$, $(v,o,\ell,r)=(28,5,4,4)$ public constants, all six maps $E_{st}$ have rank 80.  The full $80\times680$ symmetry-reduced feature map, the self-feature subspace, and the cross-feature subspace also each have rank 80.
\end{proposition}

The proposition is an exact finite-field rank certificate for the public constants.  Fix the first column $x_0$ and absorb the contributions of the remaining columns and offsets into the target.  Writing $Q_s$ for the self-quadratic contribution of column $s$, the two-column verifier fiber over the target $y$ is
\[
 \underbrace{E_{01}G_{01}(x_0,x_1)}_{\text{cross; linear in }x_1}+\;Q_1(x_1)\;=\;y-Q_0(x_0),
\]
an identity we verified directly against the official verifier on random two-column assignments.  The self-quadratic term $Q_1(x_1)$ must be retained: it is not a constant that can be moved into the target.  For fixed $x_0$ the cross term is linear in $x_1$, and in the audit its $80\times132$ coefficient matrix had rank 80, so its kernel is 52-dimensional.  On any coset of that kernel the cross term is constant, and the fiber equation restricts to the self-quadratic contribution $Q_1$ alone; the retained forms span a 50-dimensional space, giving
\[
 \boxed{50\text{ self quadrics in }52\text{ variables over }\F_{19}.}
\]
This is a restriction of the verifier fiber to a 52-dimensional affine slice, not a reparametrization of the whole fiber.  It is sound---any solution of the boxed system lifts to a genuine signature---but not complete, since the slice does not enumerate every fiber point.  The slice is underdetermined by two, so a solution is expected to exist for a generic reachable target; the value of the normal form is structural rather than cost-reducing.

The normal form is useful for both cryptanalysis and auditing because it exposes all remaining nonlinearity in a nearly square system.  It does not by itself lower the conservative headline exponent: the Hashimoto estimator for the original 50-in-102 system already specializes almost all surplus variables before its square solve.  We therefore tested whether the restricted 50-dimensional quadratic space has further public structure.  Across 100 random restricted systems, the common polar radical was always zero, the quadratic span always had dimension 50, and every individual form had rank 51 or 52.  On five independently generated keys, the adjoint/commutant algebra had dimension one.  Random sampling of $10^5$ linear combinations in the analogous Version~2.4 50-in-44 system found minimum rank 42, consistent with a generic high-rank matrix space.  These checks provide negative evidence against an immediate radical, decomposition, or low-rank shortcut, while leaving specialized multigraded solving as an open route.

\section{Concrete Complexity Estimates}\label{sec:complexity}

\subsection{Methodology}

We evaluate the new residual systems with the methodology already used in the \SNOVA{} submission and in Beullens's paper.  For a determined or overdetermined MQ system, the model estimates sparse Wiedemann/FXL from a semi-regular operating degree.  Hashimoto's method reduces an underdetermined system of $m$ equations in $n$ variables to a collection of smaller determined systems, optimizing over its parameters $(a,k)$~\cite{Hashimoto2023,Beullens2024SNOVA}.

Version~2.3 gives the Beullens cost as
\[
 \Hashimoto(n\ell-\ell^2m+R,R,q)
\]
up to its field-to-gate conversion.  We replace $(R,n\ell-\ell^2m+R)$ by the symmetry-reduced values $(K,n\ell-M+K)$ from \Cref{eq:resparams}.

The estimator admits two conventions for the semi-regular series.  The values in Version~2.3 Table~11, and the default of the public analysis script, are numerically reproduced by using no extra homogenizing variable ($p_1=0$).  Inserting one extra variable ($p_1=1$) increases the estimates by about 4.25 bits in most rows.  We adopt the stricter $p_1=1$ convention in the headline table---so our figures are conservative relative to the submission's own $p_1=0$ numbers---and give both conventions in \Cref{tab:both}.

\subsection{The determined systems actually costed by the estimator}
\label{sec:actual-cores}

The estimator does not run a Gr\"obner-basis algorithm directly on the full
underdetermined residual systems displayed in our parameter tables.  It first
specializes surplus variables and then applies HybridMQ or Hashimoto's
reduction to smaller determined or overdetermined cores.  Under the stricter
$p_1=1$ convention, the load-bearing systems are summarized in
\Cref{tab:actual-cores}.

\begin{table}[t]
\centering
\caption{Dominant determined cores selected by the current estimator under the
stricter $p_1=1$ convention.  The degree is the semi-regular operating degree
used in the cost calculation.  Degrees marked $^\dagger$ are Hashimoto blocks
that were not produced as explicit systems; they are the predicted operating
degree only, not F4-audited (\Cref{sec:validation}).}
\label{tab:actual-cores}
\small
\setlength{\tabcolsep}{4pt}
\begin{tabular}{@{}c l c c c r@{}}
\toprule
Level & Family & Public residual & Method & Dominant core & Degree\\
\midrule
I   & $\ell=4$ & $50/102$ or $50/98$ & square specialization & $50/41$ & 14\\
I   & $\ell=2$ & $48/112$ & Hashimoto $(a,k)=(2,11)$ & $46/35$ & $11^\dagger$\\
III & $\ell=4$ & $70/146$ or $70/142$ & square specialization & $70/55$ & 16\\
III & $\ell=2$ & $72/168$ & Hashimoto $(a,k)=(2,15)$ & $70/55$ & $16^\dagger$\\
V   & $\ell=4$ & $90/182$ or $90/178$ & square specialization & $90/73$ & 21\\
V   & $\ell=2$ & $96/224$ & Hashimoto $(a,k)=(2,23)$ & $94/71$ & $18^\dagger$\\
\bottomrule
\end{tabular}
\end{table}

For example, the strict Level-I $\ell=4$ estimate first specializes the $50/102$
system to a square $50/50$ system and then guesses nine more variables, leaving
the $50/41$ overdetermined core whose modeled degree is 14.  A solving-degree
experiment on the untouched $50/102$ presentation would therefore not directly
validate the quoted work factor; the relevant experiment is on these
estimator-selected cores.

\paragraph{Sensitivity to solving-degree error.}
The narrow Level-I $\ell=4$ rows are genuinely sensitive to the semi-regular
prediction.  Raising the dominant degree from 14 to 15 multiplies the squared
Macaulay dimension by approximately $\bigl(\binom{56}{15}/\binom{55}{14}\bigr)^2
=(56/15)^2$, or about $2^{3.8}$.  Thus degree 15 would place the strict estimate
close to, but still approximately below, the $2^{143}$ target; degree 16 would
likely erase the claimed shortfall.  These two rows should therefore be read as
narrow estimated shortfalls.

The larger instances have substantially more tolerance.  One additional degree
in the Level-III and Level-V dominant cores costs only roughly four bits in the
same monomial-count model, compared with margins of 22.71--44.05 bits.  The
Level-I $\ell=2$ row has a 12.70-bit margin.  Accordingly, moderate
anti-regularity could change the two narrow Level-I $\ell=4$ conclusions without
changing the Level-III or Level-V category conclusions.

\subsection{All odd-characteristic parameter sets}

For the nine $q=19$ parameter shapes in Version~2.3, \Cref{eq:resparams} gives the residual systems shown in \Cref{tab:both}.  The category targets are the values used by the submission: 143, 207, and 272 bits for Levels I, III, and V.

\begin{table}[t]
\centering
\caption{Estimates under both semi-regular conventions.  The $p_1=1$ column is the stricter convention we adopt for the headline figures and the abstract.  The $p_1=0$ column numerically matches Version~2.3 Table~11 and the default of the public analysis script.}\label{tab:both}
\small
\setlength{\tabcolsep}{4pt}
\begin{tabular}{@{}c l c r r r@{}}
\toprule
Level & Parameters & Residual MQ & $p_1=0$ & $p_1=1$ & Target\\
\midrule
I   & $(28,5,4,4)$  & $50/102$ & 134.69 & 138.94 & 143\\
I   & $(48,16,2,2)$ & $48/112$ & 126.05 & 130.30 & 143\\
I   & $(28,4,4,5)$  & $50/98$  & 134.69 & 138.94 & 143\\
III & $(40,7,4,4)$  & $70/146$ & 180.05 & 184.29 & 207\\
III & $(72,24,2,2)$ & $72/168$ & 181.89 & 184.29 & 207\\
III & $(38,5,4,5)$  & $70/142$ & 180.05 & 184.29 & 207\\
V   & $(50,9,4,4)$  & $90/182$ & 223.70 & 227.95 & 272\\
V   & $(96,32,2,2)$ & $96/224$ & 234.52 & 238.77 & 272\\
V   & $(52,6,4,6)$  & $90/178$ & 223.70 & 227.95 & 272\\
\bottomrule
\end{tabular}
\end{table}

The Level-I $\ell=4$ rows have the narrowest margin: 4.06 bits under the stricter convention.  The convention choice moves the estimate by about 4.25 bits, more than this margin, yet both conventions place these rows below target (134.69 and 138.94 bits against 143; \Cref{tab:both}), so the shortfall does not turn on that choice.  We nonetheless describe the Level-I rows as \emph{estimated below target}.  The remaining rows have margins of 12.70 to 44.05 bits under the same convention.

\subsection{Version 2.4 analysis preview}\label{sec:snova24}

The public \SNOVA{} analysis repository was updated with six parameter shapes labeled Version~2.4, although no matching specification or implementation package was public at the time of this audit~\cite{SNOVA24Preview}.  The symmetry quotient is algebraic and does not depend on the numerical parameters.  If the preview retains the symmetric public-matrix construction, the stricter $p_1=1$ convention gives \Cref{tab:snova24}.

\begin{table}[t]
\centering
\caption{Symmetry-reduced estimates for the Version~2.4 analysis preview.  These are not definitive Round~3 parameters; the table evaluates the public analysis-repository shapes under the same $p_1=1$ convention used in the headline Version~2.3 table.}\label{tab:snova24}
\small
\setlength{\tabcolsep}{5pt}
\begin{tabular}{@{}c l c r r r@{}}
\toprule
Level & Parameters $(v,o,\ell,r)$ & Residual MQ & New & Target & Margin\\
\midrule
I   & $(26,5,4,4)$   & $50/94$  & 138.94 & 143 & $-4.06$\\
I   & $(43,17,2,2)$  & $51/103$ & 141.57 & 143 & $-1.43$\\
III & $(37,7,4,4)$   & $70/134$ & 184.29 & 207 & $-22.71$\\
III & $(64,25,2,2)$  & $75/153$ & 194.57 & 207 & $-12.43$\\
V   & $(44,6,4,6)$   & $90/146$ & 227.95 & 272 & $-44.05$\\
V   & $(90,33,2,2)$  & $99/213$ & 246.84 & 272 & $-25.16$\\
\bottomrule
\end{tabular}
\end{table}

The first Level-I margin is the same narrow 4.06-bit estimate as in Version~2.3, and the second is only 1.43 bits.  Those two rows should be described as estimated shortfalls.  The Level-III and Level-V margins are substantially larger.  Reparameterizing within the same symmetric construction must therefore use $m_1\binom{\ell+1}{2}$, not $m_1\ell^2$, as the homogeneous rank ceiling.

\subsection{Interpretation of the estimates}

The estimates inherit three standard assumptions:
\begin{enumerate}
\item the residual systems behave close enough to semi-regular systems for the operating-degree calculation to be meaningful;
\item sparse linear algebra has the costs assumed by the submission's model;
\item field operations are converted to gates using the same coarse formula as the existing analysis.
\end{enumerate}
These assumptions are neither introduced nor proved by this work.  They are the methodology under which Version~2.3 claims its security levels.  Our conclusion is therefore comparative: once the homogeneous feature dimension is corrected, the candidate's own framework places the odd-characteristic sets below their categories.

The Cabarcas--Li--Verbel--Villanueva-Polanco solver exploits additional multigraded structure in \SNOVA{} systems~\cite{Cabarcas2025SNOVA}.  We do not include any further speedup from that method.

The scaled F4 anomaly reported in \Cref{app:multibase} concerns the multibase
graph-extension presentation $B_{\rm sym}(x_1,\ldots,x_c)=w+Ky$, which
introduces linearly mixed kernel variables.  It invalidates an unqualified
production-scale extrapolation for that stronger proposed solver.  It is not an
experiment on the directly eliminated one-base residual systems costed in
\Cref{tab:main,tab:both}, and therefore neither validates nor contradicts their
semi-regular model.  The two polynomial presentations are scheme-theoretically
equivalent as solution varieties, but Gr\"obner behavior need not be invariant
under such an extended graph formulation.

\section{Implementation and Algebraic Validation}\label{sec:validation}

We performed five layers of validation.

\paragraph{Universal identity.}
A self-contained checker verifies \Cref{lem:powerpair} over randomly sampled symmetric matrices and constructs the duplication map explicitly.  For $(m_1,\ell)=(5,4)$ the ordered feature dimension is 80, the unordered feature dimension is 50, and the duplication map has rank 50.

\paragraph{All-parameter quotient-rank certificates.}
For the simplest public relation ($R_0=I$ and $R_j=0$ for $j\ge1$), the quotient matrices attain full
column rank for all nine $q=19$ Version~2.3 shapes:
\[
\begin{array}{c|c|c}
(v,o,\ell,r)&\text{matrix size}&\text{rank}\\\hline
(28,5,4,4)&80\times50&50\\
(48,16,2,2)&64\times48&48\\
(28,4,4,5)&80\times50&50\\
(40,7,4,4)&112\times70&70\\
(72,24,2,2)&96\times72&72\\
(38,5,4,5)&100\times70&70\\
(50,9,4,4)&144\times90&90\\
(96,32,2,2)&128\times96&96\\
(52,6,4,6)&144\times90&90.
\end{array}
\]
For the Level-I $\ell=4$ constants, all $19^3=6{,}859$ scalar relations in the
affine chart $R_0=I$ (the standard normalization; this is a chart of
$\mathbb P^3(\F_{19})$, not all $7{,}240$ of its points) produced an
$80\times50$ matrix of rank 50, and 10,000 sampled module-valued relations did
the same.  These certificates show that the
full quotient rank used in the tables is attained, rather than extrapolated from
Level I.  Full quotient rank is the conservative case: any further rank defect
reduces the number of quadratic equations.

\paragraph{Official-key reconstruction.}
An independent Python implementation of the relevant Version~2.3 key-generation routines reproduced the Level-I $(28,5,19,4)$ KAT public key byte-for-byte.  This checks the field arithmetic, the fixed $A,B,Q$ generation (including the documented coefficient repair), and the public-matrix expansion used in the attack harness.

\paragraph{End-to-end verifier equivalence.}
On the KAT-derived key and two further generated keys, we constructed the affine-column system for a fixed public relation.  Every trial gave
\[
 \rank(\overline E_R)=50,
 \qquad \dim\ker(\overline E_R^\transpose)=30,
 \qquad \rank(W_RL_R)=30.
\]
After eliminating the 30 affine equations, a known solution satisfied the residual 50-equation system, and reconstructing the original 80 verification coordinates gave zero discrepancy.  The same harness verified that the offset construction of \Cref{lem:skew} made symmetric-block rejection impossible for every residual assignment.

\paragraph{Cross-column normal form and restricted structure.}
The complete Level-I symmetry-reduced feature map has size $80\times680$ and rank 80.  Every individual cross-column submap is $80\times80$ and invertible, yielding \Cref{prop:crosspair}.  After fixing one column, 100 independent coefficient maps all had rank 80 and a 52-dimensional kernel.  Restriction of the 50 self-quadrics to those kernels gave common-radical dimension zero in every trial, individual ranks 51 or 52, and a full 50-dimensional quadratic span.  A separate adjoint-algebra computation on five keys returned only the scalar algebra.  These experiments were added specifically to test whether the 50-in-52 normal form hides an easier decomposition; none was observed.

\paragraph{Low-degree sanity check.}
For the exported official 50-in-102 residual, the degree-one Macaulay
multiplication matrix had full column rank (5,150 columns) and no unexpected
degree-one syzygy.  More substantially, the planted residual root has full
Jacobian rank 50.  Among 10,000 random coordinate 50-dimensional strips through
that root, 9,466 had invertible restricted Jacobian, closely matching the
random-matrix probability $\prod_{j=0}^{49}(1-19^{j-50})=0.9445987429\ldots$.
Across 20 independently generated official-shaped keys, every planted root again
had Jacobian rank 50, and the number of the 10,000 sampled coordinate strips
retaining full rank ranged from 9,412 to 9,491 (mean 9,447).

All 102 coordinate directions, the planted-root direction, and 10,000 random
nonzero directions produced polar matrices of full output rank 50.  We also
evaluated every projective Hessian combination whose coefficient vector has
Hamming weight at most two, $50+\binom{50}{2}(19-1)=22{,}100$ combinations,
together with 20,000 uniformly random nonzero combinations.  The minimum
observed Hessian rank was 100 in 102 variables.  These tests provide strong
negative evidence against a common radical, a low-rank pencil, or an obvious
decomposition.

They do not by themselves prove that the production systems have the semi-regular
operating degrees in \Cref{tab:actual-cores}.  We therefore ran the direct solver
audit described above: from the official KAT keys we specialized each residual to
its estimator-selected $\ell=4$ core and compared its degree-by-degree Faug\`ere
F4 behaviour (per-degree Macaulay matrix size, new basis elements, and zero
reductions) against an ensemble of five random systems matched in field,
dimensions, quadratic support density, and consistency status.  Because a
complete solve at the modeled degree is itself of the order of the estimated
attack, no run completes; the reproducible quantity is the per-degree structure
up to the maximum degree reached under a fixed memory/time budget.

For the load-bearing Level-I $\ell=4$ core ($50/41$, modeled degree 14), the
official system is indistinguishable from the matched random ensemble through the
maximum degree reached (F4 degree~4 completed; degree~5 entered before the budget
cut both official and controls): identical degree schedule, no early degree fall,
no zero-reduction surplus beyond the F4 Koszul artifacts present identically in
the controls, and no Macaulay-rank deficit.  The official core falls within the
control ensemble's spread rather than matching any single control, consistent
with semi-regular behaviour.  This is positive evidence for the estimate, not a
solve; the row remains a narrow estimated shortfall.  The Level-III and Level-V
$\ell=4$ cores ($70/55$ and $90/73$) are larger still and lie even further beyond
the completion budget than the $50/41$ core, so no F4-audited degree is available
for them and their rows stay modeled predictions.

\paragraph{Scaling ladder.}
Because the load-bearing cores are too large to solve to completion, we validate
the operating-degree model on a ladder of \emph{faithful} symmetry-reduced cores
small enough to run msolve to a full rational parametrization.  Each core is built
through the same verified pipeline (residual construction, reduction, and
specialization), passes the identity and planted-root checks, and is kept over
$q=19$ so that field equations (active only at degree~19) never contaminate the
measured degree.  For each core we record the maximum F4 round degree and compare
against three random controls matched in dimensions, quadratic-support density,
and planted-root consistency (\Cref{tab:ladder}).

\begin{table}[t]
\centering
\caption{Observed maximum F4 round degree (msolve, run to completion) versus the
semi-regular prediction on faithful reduced cores.  $K$ is the quotient-capped
feature-map rank; each core is determined ($m=n=K$).  Daggered entries sit one
below the semi-regular first-fall index: on both, the F4 bulk table is identical
to a matching rung that reaches the predicted degree (e.g.\ the two $12/12$
cores share an identical degree schedule and peak Macaulay size), and the $\pm1$
is a trailing spair-reduction label in msolve, not a degree fall.  No core
reaches a degree \emph{above} its prediction, and each is statistically
indistinguishable from its matched random controls.}
\label{tab:ladder}
\begin{tabular}{cccccc}
\toprule
$\ell$ & $K$ ($m/n$) & density & predicted $D_\mathrm{reg}$ & observed F4 degree & random controls \\
\midrule
2 & 3 ($3/3$)    & 0.72 & 4  & 4              & 3, 4, 3   \\
2 & 6 ($6/6$)    & 0.94 & 7  & 7              & 7, 7, 7   \\
2 & 9 ($9/9$)    & 0.95 & 10 & 10             & 5, 10, 10 \\
2 & 12 ($12/12$) & 0.96 & 13 & $12^{\dagger}$ & 7, 12, 12 \\
3 & 6 ($6/6$)    & 0.90 & 7  & 7              & 4, 6, 4   \\
3 & 12 ($12/12$) & 0.95 & 13 & 13             & 7, 7, 13  \\
4 & 10 ($10/10$) & 0.96 & 11 & $10^{\dagger}$ & 6, 10, 11 \\
\bottomrule
\end{tabular}
\end{table}

On five of the seven rungs the observed solving degree equals the semi-regular
prediction exactly; on the two daggered rungs it is one lower for the bookkeeping
reason noted in the caption.  Decisively, no core solves at a degree
\emph{above} its prediction --- the signature of the Appendix-E-style
degeneration that would inflate the true attack cost.  On each rung the official
core is statistically indistinguishable from its matched controls: it lands
within the control spread on five rungs, sits one degree above the control
maximum on one ($\ell=3$, $K=6$: $7$ versus a maximum of $6$), and one below it
on one ($\ell=4$, $K=10$: $10$ versus a maximum of $11$), differences that are
within run-to-run noise.  In no case is the official core \emph{easier} than a
generic semi-regular system in a way that would cheapen the attack.  This corroborates the operating-degree model by
direct F4 measurement on genuinely reduced SNOVA cores up to $K=12$; the
extrapolation to the headline dimensions continues to rest on that (here
corroborated) predictor.

The three $\ell=2$ rows of \Cref{tab:actual-cores} are the dominant determined
\emph{blocks} of a Hashimoto change-of-basis reduction; producing those exact
blocks as explicit polynomial systems is outside the current reproducibility
tooling, so their operating degrees remain the modeled semi-regular predictions
(agreed by two independent predictors) and were \emph{not} F4-audited.  As a
structural cross-check on the same official $\ell=2$ keys we instead audited the
valid but not gate-optimal square-specialization cores ($48/36$, $72/60$,
$96/79$), which are coordinate-emittable and verified byte-for-byte against the
reference public key and are included in the reproducibility package.  These
cores are themselves above the completion budget ($n\ge36$), so they are not
F4-audited; the direct degree corroboration comes from the smaller faithful cores
of \Cref{tab:ladder}.

\section{Countermeasures and Standardization Implications}\label{sec:countermeasures}

The immediate requirement is to evaluate every symmetric-matrix parameter set using
\[
 K=m_1\binom{\ell+1}{2}
\]
as the maximum homogeneous rank in the affine-column attack.  Merely choosing $A,B,Q$ matrices for which the ordered-coordinate $E_R$ has high rank cannot overcome the quotient.

Possible design responses include:
\begin{enumerate}
\item \textbf{Reparameterization.} Increase the dimensions until the symmetry-reduced residual MQ system meets the desired category under the team's accepted cost methodology.  This is required at every level: all three NIST categories fall short, and the Level-III and Level-V shortfalls of up to 44 bits demand larger parameters than the narrow Level-I margin alone would suggest, working against the compact keys that motivated the odd-characteristic design.
\item \textbf{Remove public-matrix symmetry.} This sacrifices the key-size reduction that motivated the odd-characteristic change, but removes \Cref{lem:powerpair} in its present form.
\item \textbf{Redesign the whipping layer.} A construction with a provably full-rank emulsifier on the actual quadratic feature space, closer to the MAYO design principle, may prevent this style of dimensional collapse.  Such a redesign requires a fresh analysis; simply randomizing more $A,B,Q$ terms is insufficient.
\end{enumerate}

The symmetric-block rejection should not be counted as a countermeasure against the affine-column attack, because fixed offsets can satisfy it identically.  A rejection rule capable of blocking the attack would need to constrain the affine family itself without destroying correct signing, rather than test a property that the attacker can set with constants.

NIST advanced \SNOVA{} to Round~3 while allowing updated specifications and implementations~\cite{NISTRound3}.  The Version~2.3 draft and Version~2.4 analysis preview are therefore at an appropriate stage for correcting the attack model.  Returning unchanged to the submitted $q=16$ construction is not an obvious remedy, because the structure-aware wedge attack already reports those parameters broken.  Any viable repair must address both attack families and be analyzed on the actual, symmetry-reduced feature space.

\section{Conclusion}\label{sec:conclusion}

The public odd-characteristic \SNOVA{} redesign makes its scalar-expanded base matrices symmetric.  That key-size optimization identifies the ordered power-pair coordinates used by the Beullens affine-column attack.  The homogeneous feature map factors through
\[
 \bigoplus_{i=1}^{m_1}\Sym^2_{\F_q}(\F_q[S])
\]
and has rank at most $m_1\binom{\ell+1}{2}$, not $m_1\ell^2$.  This is a universal algebraic identity, not a weak-key phenomenon.

Under the candidate's own semi-regular Hashimoto/Wiedemann framework, every
$q=19$ Version~2.3 set and every Version~2.4 preview set is estimated below its
intended category.  The Level-III and Level-V shortfalls are robust within a
substantially wider range of solving-degree error.  The Level-I $\ell=2$ row
retains a 12.70-bit margin.  The two Level-I $\ell=4$ rows are narrower: the
strict estimator costs actual $50/41$ cores at modeled degree 14, and a rise to
degree 16 could erase their 4.06-bit shortfall.  We therefore describe those rows
specifically as narrow estimated shortfalls rather than placing all nine rows on
an equal empirical footing.

The verifier's symmetric-block rejection is bypassed by fixed skew offsets.  An
independent cross-column Gram normal form restricts the official Level-I verifier
fiber to a sound 50-quadratic system on a 52-dimensional affine slice.  Full quotient-rank
certificates hold for every Version~2.3 shape, and the official Level-I residual
exhibits full planted-root Jacobian rank, full sampled polar rank, and very high
sampled Hessian rank.  These facts rule out several obvious structural
degeneracies without proving the semi-regular operating degree.

The submitted $q=16$ package is outside this exact quotient because its public
matrices are not symmetrized, but it is already heavily affected by wedge attacks.
The present result therefore targets the public odd-characteristic repair
direction rather than offering a safe fallback.  The remaining uncertainty is
concentrated in solver cost and in the genericity assumptions the cost model rests
on.  The symmetry quotient and the rank certificates are exact, and the reduction of
a selected forgery to $K$ quadratic equations is exact when the quotient and
affine-constraint ranks are generic---the case realized in every relation we
tested---while the cross-column normal form is a sound restriction of the verifier
fiber.  The displayed gate exponents additionally inherit the heuristic
semi-regular methodology used by the \SNOVA{} submission.

\appendix
\setcounter{section}{0}
\renewcommand{\thesection}{\Alph{section}}
\renewcommand{\theHsection}{app.\Alph{section}}

\section{Parameter Formulas}\label{app:parameters}

For convenience, \Cref{tab:params} records the intermediate dimensions.  Recall
\[
 n=v+o,
 \quad m_1=\left\lceil\frac{or}{\ell}\right\rceil,
 \quad M=or\ell,
 \quad K=m_1\binom{\ell+1}{2},
 \quad N_{\rm res}=n\ell-M+K.
\]

\begin{table}[ht]
\centering
\caption{Dimension calculation for all $q=19$ Version~2.3 draft sets.}\label{tab:params}
\small
\begin{tabular}{@{}c l r r r r@{}}
\toprule
Level & $(v,o,\ell,r)$ & $m_1$ & $M$ & $K$ & $N_{\rm res}$\\
\midrule
I   & $(28,5,4,4)$  & 5  & 80  & 50 & 102\\
I   & $(48,16,2,2)$ & 16 & 64  & 48 & 112\\
I   & $(28,4,4,5)$  & 5  & 80  & 50 & 98\\
III & $(40,7,4,4)$  & 7  & 112 & 70 & 146\\
III & $(72,24,2,2)$ & 24 & 96  & 72 & 168\\
III & $(38,5,4,5)$  & 7  & 100 & 70 & 142\\
V   & $(50,9,4,4)$  & 9  & 144 & 90 & 182\\
V   & $(96,32,2,2)$ & 32 & 128 & 96 & 224\\
V   & $(52,6,4,6)$  & 9  & 144 & 90 & 178\\
\bottomrule
\end{tabular}
\end{table}

\section{A Linear-Algebra View of the Quotient}\label{app:linear}

For one base form, order the coordinates of $\F_q^{\ell^2}$ by pairs $(a,b)$.  Order the coordinates of $\Sym^2(\F_q^\ell)$ by pairs $a\le b$.  The duplication matrix $D_\ell$ has one $1$ in row $(a,a)$ and, for $a<b$, one $1$ in each of rows $(a,b)$ and $(b,a)$.  It has full column rank $\binom{\ell+1}{2}$ over every field.  For $m_1$ base forms,
\[
 D=I_{m_1}\otimes D_\ell.
\]
The attack matrix used in the odd-characteristic homogeneous system is $E_RD$.  In particular, even an ordered-coordinate matrix $E_R$ of full row and column rank cannot produce rank above $\rank(D)=K$.

For $\ell=4$ the single-form duplication matrix maps ten coordinates to sixteen positions.  The six missing dimensions are precisely the antisymmetric differences $e_{a,b}-e_{b,a}$ for $a<b$.  Across $m_1=o$ square instances, this removes $6o$ homogeneous dimensions.  For the Level-I $o=5$ set, it removes 30 dimensions, exactly the number of affine-linear equations observed in the end-to-end experiment.

\section{Why the Extension Field Does Not Give an Immediate MinRank Break}\label{app:noAshortcut}

The field notation $A=\F_q[S]\cong\F_{q^\ell}$ makes a tempting but false
inference: one might try to regard the public matrices as an $A$-linear matrix
code of size $n\times n$ and search for an $A$-rank-one dual element.  The
structural proof above does not justify that step.  The pairing
$\Phi_{u,i}$ is symmetric and bilinear over $\F_q$, but it is generally not
$A$-balanced:
\[
 \Phi_{u,i}(\alpha\xi,\eta)
 \ne
 \Phi_{u,i}(\xi,\alpha\eta).
\]
Such balance would require the public form to commute appropriately with the
$A$-action.  In \SNOVA{}, the secret transform and the $Q$ factors lie in
$A$, but the base quadratic blocks are sampled as general matrices subject
only to the UOV and, in odd characteristic, global symmetry conditions.  The
specification itself treats the extension-field decomposition as a collection
of Frobenius-conjugate lifted UOV instances, not as one $A$-linear public
matrix code.

An official-KAT audit of the Level-I $(28,5,19,4)$ key gives a particularly
sharp falsification.  With $J=I_n\otimes S$, all five public matrices satisfy
\[
 \rank(P_iJ-JP_i)=132=n\ell.
\]
None of their $5\cdot33^2=5{,}445$ scalar $4\times4$ blocks lies in the
centralizer field $A$.  The two-sided power orbit has the full tensor-square
dimension $5\ell^2=80$, while its symmetric quotient has the full dimension
$5\binom{\ell+1}{2}=50$; it does not collapse further to the
$5\ell=20$ dimensions that an $A$-balanced model would predict.  Finally, a
genuine public oil dual $uu^\transpose$ has $\F_q$-rank one, fails to commute
with $J$, and has nonzero blocks of rank one.  A nonzero multiplication matrix
from the field $A$ is invertible, so this dual is not an $A$-valued matrix.

Hence there is no justified $n\times n$ rank-one MinRank problem over
$\F_{q^\ell}$ that replaces the scalar-expanded problem.  The valid surviving
uses of the extension field are the known Frobenius/diagonal-block lifting
attacks and the $A$-valued Gram relations analyzed above.

\section{The Multibase Stable-Ideal Extension}\label{app:multibase}

\Cref{prop:cbase} gives an exact attack architecture beyond the one-base
reduction.  For a $c$-base relation, let
\[
 F_c=m_1\binom{c\ell+1}{2},\qquad M=or\ell,
 \qquad p_c=F_c-M.
\]
When $F_c\ge M$ and the public feature map is surjective, select a preimage
$w$ of the target and a basis $K$ of its kernel.  Forgery becomes
\begin{equation}\label{eq:lifted-multibase}
 B_{\rm sym}(x_1,\ldots,x_c)=w+Ky,
\end{equation}
with $p_c$ linear kernel variables.  For $c=r$, this is a symmetry-aware
stable encoding of direct forgery and requires no salt-membership search.

After diagonalizing the public $A$-action, there remain $\ell$ grading blocks,
not $c\ell$ blocks.  If each block contains $s=M/\ell=or$ specialized
quadratic variables, the numbers of feature equations of each multidegree are
\[
 A_c=m_1\binom{c+1}{2}
 \quad\text{of degree }2e_a,
 \qquad
 B_c=m_1c^2
 \quad\text{of degree }e_a+e_b.
\]
The corresponding formal complete-intersection series is
\[
 H(\mathbf t)=
 \frac{\prod_a(1-t_a^2)^{A_c}
       \prod_{a<b}(1-t_at_b)^{B_c}}
      {\prod_a(1-t_a)^s}.
\]
Using the Cabarcas-style nonhomogeneous criterion gives attractive formal
screens.  For the Level-I square $\ell=4$ set, $c=4$ gives about
$2^{99.8}$ operations with an $\omega=2$ sparse-linear-algebra model and
$2^{132.3}$ with $\omega=2.81$.  These figures are \emph{not} part of the
headline attack claim.  They depend on a solving-degree extrapolation that has
not been validated at production scale.

Scaled experiments reveal the risk.  Generic lifted systems with
$(\ell,c,m_1)=(4,4,1)$ reached maximum F4 degrees 4 and 6 when the number of
variables per grading block was one and two, respectively; the next scale
already produced very large degree-three matrices before completion.  A
separate $(3,3,1)$ family reached degree 6 where the formal criterion predicted
roughly 3--4.  These are ordinary-F4 experiments rather than a custom
multigraded implementation, so they neither prove nor refute the attack
estimate.  They do show that a public complexity claim near $2^{100}$ would be
premature without an independent multigraded-F4 or Magma validation.

The graph extension is scheme-theoretically equivalent to direct forgery:
projection onto the $x$-coordinates identifies its solution variety with the
corresponding fiber of the direct feature map.  It should not, however, be
confused with the directly eliminated one-base polynomial presentation.  The
public kernel generally mixes the nominal grading sectors, so the graph-extension
variables do not decompose into the independent homogeneous blocks assumed by the
most optimistic formal complete-intersection screen.  Gr\"obner behavior can
therefore differ substantially even though the two presentations encode the same
solutions.

The exact conclusions are therefore: the multibase factorization, feature counts,
graph equivalence, and surjectivity tests are valid; the scaled F4 anomaly applies
to this extended formulation; and the stronger multibase work factors remain an
unconfirmed solver extrapolation.  It supplies no direct evidence for or against
the one-base estimator cores in \Cref{tab:actual-cores}.

\section{Reproducibility and Claim Checklist}\label{app:repro}

A public claim based on this work should distinguish the following statements.

\begin{description}[leftmargin=0pt]
\item[Proved algebraically.] The homogeneous affine-column system factors through unordered power pairs, so its rank is at most $K$.  The verifier's symmetric-block rejection can be bypassed by fixed skew offsets.
\item[Checked computationally.] The Level-I official-key reconstruction matches a KAT public key; three end-to-end trials yield an exact 50-in-102 system; all six official Level-I cross-column blocks are invertible and yield the 50-in-52 normal form; all nine Version~2.3 shapes have full quotient-rank certificates; and the official Level-I residual has full sampled polar rank, full planted-root Jacobian rank, and high-rank sampled Hessian combinations.
\item[Estimated heuristically.] The gate complexities in \Cref{tab:main,tab:both}, like the submission's estimates, assume semi-regular behavior and a particular sparse-linear-algebra model.  The two Level-I $\ell=4$ rows are sensitive to a one- or two-degree error; the Level-III and Level-V margins tolerate several additional degrees in the same model.
\item[Audited negative result.] The public forms are not $A$-linear and there is no immediate $n\times n$ MinRank shortcut over $\F_{q^\ell}$.
\item[Preliminary extension.] The exact multibase Gram attack is included as a solver-research direction, but its stronger formal exponents are not claimed as validated production work factors.
\item[Not claimed.] We do not claim a practical production-size forgery, a key-recovery attack, a break of the $q=16$ parameter sets, or a proven running-time upper bound matching the estimates.
\end{description}

The accompanying reproducibility package contains:
\begin{itemize}
\item a self-contained proof-of-concept checker for the symmetry identity, quotient map, parameter formulas, and rejection bypass;
\item a clean implementation of both semi-regular conventions and the Hashimoto optimization;
\item exact cost tables;
\item rank-search logs;
\item the low-degree sanity-check script (Jacobian, coordinate-strip, polar, and Hessian statistics, plus the degree-one Macaulay column-rank check) and its logs;
\item the solving-degree experiment: official $\ell=4$ cores specialized from the KAT keys, matched random-control ensembles, and the per-degree F4 round tables with the maximum degree reached under budget;
\item end-to-end KAT and verifier-equivalence transcripts;
\item SHA-256 hashes identifying the Version~2.3 draft and evidence files used in the audit.
\end{itemize}
The clean official-source reimplementation will be released with the public version after coordinated disclosure.

The SHA-256 digest of the Version~2.3 PDF used in the implementation audit is
\begin{center}
\ttfamily\small 963fa43d7e6ed8ecf1e77556e7cf47f62d4c34177efece2cd9a9517d0fbaf783.
\end{center}

\printbibliography

\end{document}
